\section{Conclusion}
\label{sec:conclusion}
La question initiale ne reçoit pas une réponse de supériorité uniforme.
Dans les cas synthétiques étudiés, la dépendance des échanges aux
variations locales modifie le compromis entre bruit et structures,
mais cet effet dépend fortement de $K$, de la conductance et du temps.
À paramètres fixés, la loi rationnelle avec $K=0{,}10$ et $n=20$
sur l'image géométrique bruitée à $\sigma=0{,}10$ améliore les deux
mesures principales par rapport à la chaleur. La loi exponentielle
au même réglage améliore le PSNR mais diminue le SSIM relatif à la
chaleur. Ce désaccord montre pourquoi le jugement ne peut être réduit
à une unique métrique.
Les faibles conductances peuvent protéger des différences dues au
bruit ; une diffusion trop longue peut au contraire effacer des
transitions. Le profil et les cartes illustrent ces mécanismes sans
constituer une nouvelle mesure quantitative de contours. Les meilleurs
réglages de la grille ont une valeur exploratoire : leur sélection
utilise l'image propre et certains maxima atteignent la frontière
temporelle étudiée. Ils ne définissent donc ni optimum global ni
règle de déploiement.

Sur le plan numérique, le partage des flux garantit la conservation,
et la borne sur le pas explicite donne le principe du maximum discret.
Ces propriétés sont établies pour le schéma réellement exécuté et
confortées par les tests existants. Elles ne prouvent pas la bonne
position de l'EDP de Perona--Malik ni une approximation isotrope de
son gradient : l'objet évalué est une variante directionnelle.
Le résultat principal de ce projet est ainsi une comparaison
conditionnelle et traçable, articulant modèle, démonstration discrète,
vérification et interprétation. L'extension à des images naturelles
et le choix des paramètres sans référence propre restent des
perspectives, au-delà du périmètre synthétique étudié.
